% Called after \appendix, which prints the title and appendix contents page.

\section{Additional methods}
\label{app:methods}
This appendix demonstrates how supporting material follows the main article.
Replace these instructional paragraphs with experimental details, derivations
and supplementary analyses. The typography, citation colour and paragraph
indentation follow the main text; page numbers continue without restarting.

\subsection{Experimental protocol}
\label{app:protocol}
Describe the preparation of materials, the sampling procedure and any criteria
used to exclude observations. Report the order of operations and the controls
needed to reproduce the study. If a procedure is adapted from earlier work,
state which steps were changed and why.

Provide the details that would interrupt the main narrative but remain necessary
for reproducibility. These may include instrument settings, preprocessing rules,
software versions and parameter choices. Separate prespecified decisions from
exploratory analyses, and explain how missing observations were handled.

\subsection{Reporting checklist}
Table~\ref{tab:supplement} illustrates the layout of a supplementary table.
Table and equation numbers continue from the main article in this template.
No extra numbering commands are required.

\begin{table}[htbp]
  \centering\small
  \caption{Supplementary reporting checklist. This is a layout example,
    not a record of experimental results.}
  \label{tab:supplement}
  \begin{tabular*}{\linewidth}{@{\extracolsep{\fill}}ll@{}}
    \toprule
    Item & Information to provide \\
    \midrule
    Data & Source and version \\
    Sampling & Inclusion criteria \\
    Analysis & Parameters and software \\
    Repetition & Seeds and uncertainty \\
    \bottomrule
  \end{tabular*}
\end{table}

\section{Supporting derivation}
\label{app:derivation}
Starting from \Eqref{eq:objective}, the gradient of the empirical objective
is, when each summand is differentiable,
\begin{equation}
  \nabla_{\theta}\mathcal{L}(\theta)
  =\frac{1}{N}\sum_{i=1}^{N}
  \nabla_{\theta}\ell\bigl(f_{\theta}(x_i),y_i\bigr).
  \label{eq:supplement-gradient}
\end{equation}
This identity follows by differentiating the finite sum term by term. State any
additional assumptions before using an analogous step in your own derivation,
and explain how the result supports the claims in the main article.

Appendix~\ref{app:methods} and \Eqref{eq:supplement-gradient} demonstrate
linked references within the supporting material. Use the same notation as in
the main text so that readers can move between the argument and its details.
